\documentclass[12pt]{article}
\usepackage{amsmath, amssymb, amsthm}
\usepackage{bm}
\usepackage{graphicx}
\usepackage{geometry}
\geometry{a4paper, margin=1in}
\usepackage{natbib}
\usepackage{hyperref}

\title{A Matrix-Free Mixed-Frequency Approach to \\ Order-Invariant SVMVARs}
\author{Independent Research Report}
\date{\today}

\begin{document}
\maketitle

\begin{abstract}
This report develops a methodological research agenda for extending the Order-Invariant Stochastic Volatility in Mean Vector Autoregression (OI-SVMVAR) of \citet{davidson2025investigating} to a mixed-frequency setting. To reduce temporal aggregation bias, the model is written as a high-frequency latent state-space system. The proposed implementation combines the matrix-valued basis-lag representation of \citet{zhu2026structural} with Gaussian latent-state samplers \citep{papandreou2010gaussian,bardsley2012mcmc}. The extension is not a mechanical application of Zhu's sampler: Zhu's baseline is homoskedastic and explicitly leaves stochastic volatility as future work. The key computational idea is that SV destroys Zhu's time-invariant middle precision block, but not the matrix-free factorisation itself: the operator becomes a time-invariant VAR shell with a time-varying sequence of small precision blocks. Pilot benchmarks show that this observation is not sufficient for a speed claim. Lightweight PCG preconditioners have iteration counts that explode with SV dispersion, while direct sparse Cholesky is the practical production solver at Taiwan-scale monthly dimensions. Matrix-free PCG remains a memory and larger-scale fallback, conditional on stronger global preconditioning.
\end{abstract}

\section{Introduction}
The \citet{davidson2025investigating} OI-SVMVAR model effectively disentangles the real effects of macroeconomic and financial uncertainty without imposing arbitrary recursive orderings. However, its single-frequency nature forces researchers to aggregate timely financial data to a quarterly clock. As \citet{zhu2026structural} demonstrates, temporal aggregation can obscure immediate financial responses and distort structural identification. 

This report sets out a Matrix-Free MF-OI-SVMVAR and records the extra obligations created by combining these literatures. The high-frequency representation preserves the timing information that is lost under low-frequency aggregation, but it does not automatically inherit all identification and computational results from the source papers. In particular, \citet{davidson2025investigating} prove point identification in an observed single-frequency system, while \citet{zhu2026structural} develop the matrix-free mixed-frequency sampler under homoskedastic innovations. The computational contribution is narrower than originally hoped: the SV-aware matrix-free product exists, but simple preconditioners do not make PCG competitive at the target Taiwan-scale dimensions. The practical near-term implementation should therefore benchmark direct sparse Cholesky as the production latent-state solver and reserve matrix-free PCG for memory-constrained or larger-frequency-gap regimes. The identification contribution is harder and remains separate: one must show that DHK's heteroskedasticity-based identification survives the latent mixed-frequency likelihood rather than only the imputed high-frequency conditional likelihood.

\section{The High-Frequency Latent State-Space}
We assume the true economic dynamics operate at a high frequency (e.g., monthly). Let $\mathbf{y}_t^*$ be an $n \times 1$ vector of latent high-frequency variables, partitioned into fully observed high-frequency financial variables and partially observed low-frequency macroeconomic variables (e.g., GDP). The structural model follows the OI-SVMVAR process:
\begin{equation}
    \mathbf{y}_t^* = \mathbf{b}_0 + \sum_{i=1}^p \mathbf{B}_i \mathbf{y}_{t-i}^* + \sum_{j=0}^q \mathbf{A}_j \mathbf{h}_{t-j} + \mathbf{B}_0^{-1} \bm{\epsilon}^y_t, \quad \bm{\epsilon}^y_t \sim \mathcal{N}(\mathbf{0}, \mathbf{U}_t)
\end{equation}
where $\mathbf{B}_0$ is the order-invariant contemporaneous impact matrix with ones on the diagonal and unconstrained off-diagonal elements. The intercept $\mathbf{b}_0$ is added as a location term relative to the baseline DHK exposition; it can be absorbed into the stacked intercept without affecting the identification argument conditional on the latent path.

Stacking the observations over time $t=1,\dots,T$, we can express the system in matrix form:
\begin{equation}
    \mathbf{H}_{\mathbf{B}} \mathbf{y}^* = \mathbf{A}_{\mathbf{h}} \mathbf{h} + \mathbf{c}_{\mathbf{B}} + (\mathbf{I}_T \otimes \mathbf{B}_0^{-1}) \bm{\epsilon}^y, \quad \bm{\epsilon}^y \sim \mathcal{N}(\mathbf{0}, \mathbf{U})
\end{equation}
where $\mathbf{H}_{\mathbf{B}}$ is a banded $Tn \times Tn$ matrix built from the VAR coefficient matrices $\mathbf{B}_i$, $\mathbf{c}_{\mathbf{B}}$ contains intercepts and initial conditions, and $\mathbf{U}=\operatorname{diag}(\mathbf{U}_1,\ldots,\mathbf{U}_T)$ is the full block-diagonal stochastic-volatility covariance. Equivalently, the innovation precision used below is
\begin{equation}
    \mathbf{D} = \operatorname{diag}(\mathbf{D}_1,\ldots,\mathbf{D}_T), 
    \qquad \mathbf{D}_t = \mathbf{B}_0' \mathbf{U}_t^{-1} \mathbf{B}_0 .
\end{equation}
This notation is deliberately not $\mathbf{I}_T\otimes \mathbf{D}$ unless the shocks are homoskedastic.

To distinguish directly observed and missing high-frequency entries, write
\begin{equation}
    \mathbf{y}^*=\mathbf{S}^o\mathbf{y}^o+\mathbf{S}^m\mathbf{y}^m ,
\end{equation}
where $\mathbf{y}^o$ collects entries observed at the high-frequency clock, $\mathbf{y}^m$ collects latent entries, and $\mathbf{S}^o,\mathbf{S}^m$ are selection matrices. Following \citet{mariano2003new} and \citet{zhu2026structural}, low-frequency observations enter through a sparse aggregation matrix. After substituting the observed high-frequency entries, the low-frequency measurement equation can be written as
\begin{equation}
    \tilde{\mathbf{y}}^L = \mathbf{M}_m \mathbf{y}^m + \mathbf{v}, 
    \quad \tilde{\mathbf{y}}^L \equiv \mathbf{y}^L-\mathbf{M}_o\mathbf{y}^o,
    \quad \mathbf{v} \sim \mathcal{N}(\mathbf{0}, \lambda^{-1} \mathbf{I}) .
\end{equation}
For flow variables in growth rates, rows of $\mathbf{M}_m$ contain the Mariano--Murasawa triangular weights; for stock variables they are selection rows. The finite-$\lambda$ version is best interpreted as a numerical penalty around an exact aggregation restriction, not as ordinary measurement error. The limit $\lambda\rightarrow\infty$ recovers the exact aggregation hyperplane, while large finite $\lambda$ may worsen conditioning and must be checked by sensitivity analysis.

\section{Statistical Bottleneck: Parametric Lag Expansion}
A monthly VAR requires a large lag order $p$ to capture sluggish macroeconomic propagation. To prevent parameter proliferation ($p n^2$ free parameters), we represent the lag sequence $\mathbf{B}_1, \dots, \mathbf{B}_p$ as a linear combination of a small number of known basis functions evaluated on a normalised grid:
\begin{equation}
    \mathbf{B}_j = \sum_{k=1}^K \mathbf{\Gamma}_k \phi_k(j/p), \quad j = 1, \dots, p
\end{equation}
where $\phi_k(\cdot)$ are known basis functions and $\mathbf{\Gamma}_k$ are $n \times n$ coefficient matrices. The matrix-valued VAR version follows \citet[][Section 3.1]{zhu2026structural}, who motivates the sieve representation with approximation results such as \citet{chen2007large}. The related TVP-MIDAS paper of \citet{chan2025time} supports the broader use of linear-in-parameters basis expansions, but it is a scalar single-equation MIDAS framework with Fourier and Almon bases rather than a VAR with matrix-valued lag coefficients.

This representation replaces $p n^2$ autoregressive parameters with $K n^2$ basis parameters. It does not by itself guarantee stationarity. The posterior simulation must either impose a stationarity-preserving prior on $\{\mathbf{\Gamma}_k\}_{k=1}^K$ or reject draws for which the companion matrix implied by $\{\mathbf{B}_j\}_{j=1}^p$ has spectral radius at least one. Near-unit-root draws are also computationally costly because they increase the condition number of the latent-state precision and hence the number of PCG iterations.

\section{Matrix-Free State Sampling}
In conventional mixed-frequency settings, drawing the latent high-frequency path involves a precision matrix. Using the selection notation above, let
\begin{equation}
    \mathbf{A}=\mathbf{H}_{\mathbf{B}}\mathbf{S}^m,
    \qquad 
    \mathbf{r}=\mathbf{A}_{\mathbf{h}}\mathbf{h}+\mathbf{c}_{\mathbf{B}}-\mathbf{H}_{\mathbf{B}}\mathbf{S}^o\mathbf{y}^o .
\end{equation}
The penalised precision for the missing path is
\begin{equation}
    \overline{\mathbf{K}} = \mathbf{A}'\mathbf{D}\mathbf{A} + \lambda \mathbf{M}_m'\mathbf{M}_m .
\end{equation}
When bridging large frequency gaps with long lags, $\overline{\mathbf{K}}$ becomes wide-band. Direct sparse Cholesky factorisation \citep{chan2009efficient} scales poorly with the bandwidth square.

\subsection{SV-Aware Matrix-Vector Product}
Following \citet{zhu2026structural}, direct factorisation can be replaced by a matrix-free precision product. The important point is that stochastic volatility changes the middle precision block, not the VAR shell. The operator can be applied as
\begin{enumerate}
    \item Embed a trial vector $\mathbf{x}$ into the missing high-frequency locations and apply the forward VAR operator: $\mathbf{w}=\mathbf{A}\mathbf{x}$.
    \item For each high-frequency date, apply the local SV precision block: $\mathbf{u}_t=\mathbf{D}_t\mathbf{w}_t$, where $\mathbf{D}_t=\mathbf{B}_0'\mathbf{U}_t^{-1}\mathbf{B}_0$.
    \item Apply the adjoint VAR operator: $\mathbf{z}=\mathbf{A}'\mathbf{u}$.
    \item Add the soft-aggregation contribution: $\overline{\mathbf{K}}\mathbf{x}=\mathbf{z}+\lambda\mathbf{M}_m'\mathbf{M}_m\mathbf{x}$.
\end{enumerate}
This product remains matrix-free because no $Tn\times Tn$ precision matrix is formed. The extra SV cost is the local multiplication by $T$ small $n\times n$ matrices, which is $\mathcal{O}(Tn^2)$ under dense contemporaneous covariance blocks. For moderate $n$, this is not a qualitative change relative to the forward and adjoint basis-filter operations. It is, however, a substantive extension of Zhu: he notes that the time-invariant precision operator in his homoskedastic baseline would not survive a daily stochastic-volatility process. In the present proposal, that warning is not a dead end; it defines the open algorithmic subproblem.

\subsection{Preconditioning and Complexity}
The computational claim must therefore be stated conditionally. Zhu's best case is $\mathcal{O}(Tn)$ memory and $\mathcal{O}(\text{iter}_{PCG}KT(\log T+n))$ work per draw when basis filters are implemented with FFTs. In the present SV-in-mean extension, a more explicit accounting adds the local block cost, giving\footnote{The $\log T$ term assumes Zhu's FFT-based basis-filter implementation. The current pilot validation uses a materialised sparse $\mathbf{H}_{\mathbf{B}}$ shell and does not by itself realise this term.}
\begin{equation}
    \mathcal{O}\left(\text{iter}_{PCG}\{KT(\log T+n)+Tn^2+\operatorname{nnz}(\mathbf{M}_m)\}\right)
\end{equation}
per draw. Thus this complexity expression is a proposal-facing target rather than an already validated cost for the present code path. The term that can destroy the speedup is not the existence of the $Tn^2$ local block multiplication but the behaviour of $\text{iter}_{PCG}$.

Pilot benchmarks on the Gaussian latent-state problem confirm this risk. With $n=5$, $T=480$, $p=2$, $\lambda=10^4$, and target VAR radius 0.9, unpreconditioned PCG requires roughly 2500 iterations at moderate SV dispersion and reaches the iteration cap under high dispersion. Scalar Jacobi and date-wise block-Jacobi improve this only marginally. A time-averaged sparse-factor preconditioner reduces iterations substantially at moderate dispersion, but at Taiwan-scale monthly dimensions it is still slower than direct sparse Cholesky. The immediate production recommendation is therefore direct sparse Cholesky for the latent-state Gaussian solve, with the matrix-free operator retained for memory diagnostics, assembly-cost benchmarks, and larger-scale fallback experiments.

A natural lightweight preconditioner is block-Jacobi,
\begin{equation}
    \mathbf{P}_{BJ}=\operatorname{blockdiag}(\overline{\mathbf{K}}_{11},\ldots,\overline{\mathbf{K}}_{TT}),
\end{equation}
where each diagonal block is computed from the local contribution of $\mathbf{D}_t$ and the local aggregation penalty. The benchmark evidence indicates that this local scaling is not enough: the hard conditioning comes from the global lag-shell structure in $\mathbf{A}'\mathbf{D}\mathbf{A}$, not only from date-specific volatility scales. Any future matrix-free speed claim therefore requires a stronger preconditioner, such as an approximate sparse factor built from a time-averaged or otherwise smoothed precision operator, and must be evaluated against a direct Cholesky baseline that includes explicit $\overline{\mathbf{K}}$ assembly cost.

\section{High-Frequency Identification of \texorpdfstring{$\mathbf{B}_0$}{B0}}
Conditional on a sampled high-frequency path $\mathbf{y}^*$, the DHK row-by-row sampler can be applied exactly as in a single-frequency OI-SVMVAR. The identification of $\mathbf{B}_0$ then relies on heteroskedasticity in the structural shocks \citep{lutkepohl2020bayesian,davidson2025investigating}: the diagonal volatility paths must provide enough independent variation to satisfy the usual rank condition, with the remaining sign and permutation normalisations fixed by the unit diagonal and labelling convention.

The unconditional mixed-frequency identification claim is weaker. DHK's proof is for observed monthly data, whereas the present likelihood integrates over latent high-frequency macro paths subject to aggregation restrictions. Thus the model should be described as preserving the DHK identification argument \emph{conditional on} $\mathbf{y}^*$ and as requiring an additional mixed-frequency identification check. Two diagnostics are necessary: first, verify the heteroskedastic rank condition across posterior draws of $\mathbf{y}^*$; second, examine whether the condition remains stable as $\lambda\rightarrow\infty$, when the latent path is restricted to the exact Mariano--Murasawa aggregation hyperplane. This moves the proposal away from the over-strong claim that high-frequency modelling automatically eliminates the low-frequency identification conflict.

\section{Posterior Simulation: The MCMC Algorithm}
Combining the high-frequency state-space with the SVMVAR structure alters the joint likelihood, requiring a custom block-Gibbs sampling scheme. Conditional on the structural parameters, the model remains a linear Gaussian state-space, but the mean equations are augmented by the stochastic volatility components. The MCMC algorithm cycles through the following steps:

\subsection{Step 1: Drawing the Latent High-Frequency Path \texorpdfstring{$(\mathbf{y}^*)$}{(y*)}}
Conditional on all parameters and volatilities, the missing high-frequency path satisfies the Gaussian system
\begin{equation}
    \mathbf{A}\mathbf{y}^m = \mathbf{r}+\bm{\eta}_1, 
    \quad \bm{\eta}_1\sim\mathcal{N}(\mathbf{0},\mathbf{D}^{-1}),
    \qquad
    \tilde{\mathbf{y}}^L = \mathbf{M}_m\mathbf{y}^m+\bm{\eta}_2,
    \quad \bm{\eta}_2\sim\mathcal{N}(\mathbf{0},\lambda^{-1}\mathbf{I}).
\end{equation}
The conditional posterior is $\mathbf{y}^m \sim \mathcal{N}(\bm{\mu}_{m}, \overline{\mathbf{K}}^{-1})$, where $\overline{\mathbf{K}}=\mathbf{A}'\mathbf{D}\mathbf{A}+\lambda\mathbf{M}_m'\mathbf{M}_m$. The mean vector satisfies
\begin{equation}
    \overline{\mathbf{K}} \bm{\mu}_{m} = \mathbf{A}' \mathbf{D} \mathbf{r} + \lambda \mathbf{M}_m' \tilde{\mathbf{y}}^L .
\end{equation}
For the perturbation-optimisation draw \citep{papandreou2010gaussian,bardsley2012mcmc}, sample independent $\bm{\eta}_1$ and $\bm{\eta}_2$ from the distributions above and solve
\begin{equation}
    \overline{\mathbf{K}} \mathbf{y}^{m,draw}
    =
    \mathbf{A}'\mathbf{D}(\mathbf{r}+\bm{\eta}_1)
    +\lambda \mathbf{M}_m'(\tilde{\mathbf{y}}^L+\bm{\eta}_2)
\end{equation}
by PCG using the matrix-free product described in Section 4. The full latent path is then reconstructed as $\mathbf{y}^{*,draw}=\mathbf{S}^o\mathbf{y}^o+\mathbf{S}^m\mathbf{y}^{m,draw}$.

\subsection{Step 2: Drawing the Order-Invariant Impact Matrix \texorpdfstring{$(\mathbf{B}_0)$}{(B0)}}
Conditional on the latent high-frequency path $\mathbf{y}^*$, we define the SVMVAR residuals $\hat{\mathbf{y}}_t = \mathbf{y}_t^* - \sum_{i=1}^p \mathbf{B}_i \mathbf{y}_{t-i}^* - \sum_{j=0}^q \mathbf{A}_j \mathbf{h}_{t-j}$. We then apply the parameter transformation and mixture approximation algorithm of \citet{davidson2025investigating} row-by-row to $\mathbf{B}_0$, treating $\hat{\mathbf{y}}_t$ as fully observed high-frequency data.

\subsection{Step 3: Drawing Parametric Lag Coefficients \texorpdfstring{$(\mathbf{\Gamma}_k)$}{(Gamma k)}}
Conditional on $\mathbf{y}^*$, $\mathbf{h}$, and $\mathbf{B}_0$, the system is linear in the basis coefficients $\mathbf{\Gamma}_k$. Define the basis-weighted lag regressors
\begin{equation}
    \mathbf{z}_{k,t}=\sum_{j=1}^p \phi_k(j/p)\mathbf{y}_{t-j}^*,
    \qquad k=1,\ldots,K,
\end{equation}
and subtract the volatility-in-mean term,
\begin{equation}
    \mathbf{y}_t^\dagger
    =
    \mathbf{y}_t^*-\mathbf{b}_0-\sum_{j=0}^q\mathbf{A}_j\mathbf{h}_{t-j}.
\end{equation}
Then
\begin{equation}
    \mathbf{y}_t^\dagger
    =
    \sum_{k=1}^K \mathbf{\Gamma}_k\mathbf{z}_{k,t}
    +
    \mathbf{B}_0^{-1}\bm{\epsilon}^y_t,
    \qquad \bm{\epsilon}^y_t\sim\mathcal{N}(\mathbf{0},\mathbf{U}_t).
\end{equation}
This is a SUR update with time-varying covariance after conditioning on $\mathbf{B}_0$ and $\{\mathbf{U}_t\}$, equipped with shrinkage priors such as Minnesota or Horseshoe priors. Stationarity screening is applied after reconstructing $\mathbf{B}_1,\ldots,\mathbf{B}_p$ from the proposed $\mathbf{\Gamma}_k$ draw.

\subsection{Step 4: Drawing Common Volatilities \texorpdfstring{$(\mathbf{h})$}{(h)} and Classifications}
Conditional on $\mathbf{y}^*$ and the VAR parameters, the extraction of the common macro and financial volatilities $\mathbf{h}_t$, the idiosyncratic volatilities, and the Markov-switching classification states follows exactly the band-sparse Acceptance-Rejection Metropolis-Hastings steps developed by \citet{davidson2025investigating}, operated entirely at the high frequency.

\section{Methodological Strategy: Finite-Sample Inference}
Following the established paradigm of large Bayesian VAR literature, our approach explicitly departs from large-sample asymptotics:
\begin{enumerate}
    \item \textbf{Guaranteed Posterior Propriety via Proper Priors:} Because $\mathbf{B}_0$ is unrestricted, we circumvent divergent posteriors by strictly specifying \textit{proper priors} (e.g., Gaussian, Horseshoe) for all parameters, mathematically guaranteeing a proper posterior distribution.
    \item \textbf{Focus on Finite-Sample Regularization:} In high-dimensional regimes, traditional $T \to \infty$ asymptotics fail. Our methodology embraces Bayesian shrinkage and parametric lag restriction as finite-sample regularization techniques.
    \item \textbf{Validation via Simulation:} The validity and computational tractability of the MF-OI-SVMVAR must be demonstrated empirically through Monte Carlo simulations. The design should vary SV dispersion directly and report PCG iteration counts, total wall-clock time against direct sparse Cholesky including assembly cost, memory use, sensitivity to $\lambda$, stability rejection rates, and the posterior rank diagnostics for $\mathbf{B}_0$.
\end{enumerate}

\section{Conclusion and Next Steps}
This document is a revised theoretical blueprint rather than a completed proof. It identifies a feasible route to an MF-OI-SVMVAR, but the hard parts are precisely those raised by the source papers. On the computational side, the task is not to rescue matrix-free sampling from stochastic volatility; the matrix-vector product still exists. The current evidence says that lightweight PCG is not the production path at Taiwan-scale monthly dimensions. Direct sparse Cholesky should be the default solver unless a fair total-cost and memory benchmark shows that assembly or memory pressure reverses the ranking. Matrix-free PCG remains useful as a larger-scale fallback and as a preconditioner-design research problem. On the identification side, the harder problem is proving or diagnosing mixed-frequency identification of $\mathbf{B}_0$, because DHK's proof applies at the observed frequency while this likelihood integrates over latent high-frequency macro paths.

\bibliographystyle{aea.bst}
\bibliography{references}

\end{document}
